\documentclass[11pt,a4paper]{article}

% ---------- Packages ----------
\usepackage[utf8]{inputenc}
\usepackage[T1]{fontenc}
\usepackage{lmodern}
\usepackage[margin=1in]{geometry}
\usepackage{titlesec}
\usepackage{enumitem}
\usepackage{xcolor}
\usepackage{hyperref}
\usepackage{parskip}
\usepackage{microtype}

% ---------- Colours ----------
\definecolor{accent}{RGB}{34,58,102}

% ---------- Hyperlinks ----------
\hypersetup{
  colorlinks=true,
  linkcolor=accent,
  urlcolor=accent,
  citecolor=accent
}
\urlstyle{same}

% ---------- Section styling ----------
\titleformat{\section}
  {\color{accent}\large\bfseries}{\thesection.}{0.6em}{}
\titlespacing*{\section}{0pt}{1.4em}{0.5em}
\titleformat{\subsection}
  {\color{accent}\normalsize\bfseries}{\thesubsection}{0.6em}{}
\titlespacing*{\subsection}{0pt}{1.0em}{0.35em}

\setlist[itemize]{leftmargin=1.4em,itemsep=0.25em,topsep=0.35em}

\begin{document}

% ---------- Header block ----------
\begin{center}

{\color{accent}\fontsize{36}{40}\selectfont\textbf{Draftly}}\\[0.55em]

{\color{accent}\large\itshape
A Lawyer-in-the-Loop Legal Document Drafting and \\[0.15em]
Review Platform for Sri Lankan Legal Workflows
}\\[1.0em]

{\normalsize
\textbf{Group 06}\\[0.3em]
Dhanapala D.H.N., Dilshan A.D.L., De Silva B.K.P.
}\\[0.1em]

{\normalsize
\textbf{Supervisor:} Dr. Nisansa de Silva\\
\textbf{Teaching Assistant:} Ovindu Atukorala\\[1em]
Department of Computer Science and Engineering\\
University of Moratuwa, Sri Lanka
}\\[0.8em]

\end{center}

\vspace{0.7em}
{\color{accent}\rule{\linewidth}{1pt}}

% =====================================================
\section{Introduction}

In Sri Lanka, many legal workflows still depend on repeated manual drafting and review.
Lawyers and legal clerks prepare documents such as deeds, title reports, affidavits,
petitions, legal notices, certificates and letters by reading supporting material,
extracting the relevant facts, checking them, and then adapting previous drafts or
templates. The difficult part is not only typing the final document; it is identifying
the correct facts from client-specific records and ensuring that the draft is complete,
consistent and ready for professional review.

Over the last few years, several Sri Lankan legal platforms have appeared, such as
LawLanka, AI Pazz, Paralegal.lk and
LawCeylon~\cite{lawlanka,aipazz,paralegal,lawceylon}.
These primarily support legal research,
case law access and legislation retrieval, and they are widely used in practice.
Draftly is designed to complement these tools rather than replace them: it focuses on
client-specific legal document drafting and review, where the required content must be
derived from the user's own uploaded documents and checked by a lawyer.

For the first use case, we focus on conveyancing and title-documentation workflows.
This is a suitable starting point because lawyers must derive outputs such as title
reports, pedigrees, abstracts of title, ownership certificates and assessor letters from
prior deeds, survey plans and related records. It also gives the project a clear data
science component, because facts must be extracted and linked across documents before
any draft can be generated.

To confirm that this gap is real, we have conducted preliminary discussions with
practicing Sri Lankan lawyers, notaries and legal academics specialising in
conveyancing. Their feedback indicates that while legal research platforms are useful,
there remains a practical gap in preparing repeated legal documents from
client-specific records. This is the gap Draftly is built to explore.

\vspace{0.3em}
\noindent\emph{In short: Draftly is a lawyer-in-the-loop drafting and review platform
that turns uploaded legal documents and verified facts into reviewable drafts. In the
first use case, it applies this idea to conveyancing and title documentation.}

% =====================================================
\section{Proposed Solution}

Draftly is a \emph{lawyer-in-the-loop} platform: the system proposes and the lawyer
reviews, edits and approves. No output is ever treated as a final legal document on its
own; every draft is presented for a qualified professional to check and sign off on.

In a typical workflow, the lawyer or legal clerk creates a matter, uploads supporting
documents, reviews the facts suggested by the system, corrects them where necessary and
then generates an editable draft from an approved template. For the initial
conveyancing use case, the supporting documents may include previous deeds, survey
plans and assessment records, and the outputs may include a title report, pedigree,
abstract of title, ownership certificate or assessor letter.

\subsection{Research Novelty}
The template is only the final output layer. The more important part is the earlier
document-understanding workflow: reading uploaded material, extracting structured
facts, linking related facts across documents, identifying missing or inconsistent
information, and showing the lawyer where each fact came from.

This makes Draftly different from a simple form-filling system. In a form-filling
system, the user manually types all values into text boxes and the system places them in
a document. In Draftly, the system attempts to derive candidate facts from uploaded
documents, presents them with evidence, and lets the lawyer verify or correct them
before a draft is produced. In the conveyancing use case, this includes tasks such as
extracting deed numbers, parties, dates, plan numbers, lot numbers and extents, then
using those facts to support title reports, pedigrees and related documents.

\subsection{Key Features}
\begin{itemize}[itemsep=0pt, topsep=0pt]
  \item Upload and organise legal documents for a matter
  \item Extract important facts from scanned and digital legal documents
  \item Show extracted facts with source evidence for lawyer verification
  \item Highlight missing or inconsistent information
  \item Generate editable drafts from approved templates and verified facts
  \item Support lawyer review, correction, approval and export
\end{itemize}

\subsection{Data Science Component}
The data science component focuses on document understanding and fact extraction from
Sri Lankan legal documents. For the first use case, this includes OCR for scanned deeds
and forms, extraction of parties, deed numbers, dates, notaries, property details,
registration references and survey-plan details, and linking those facts across a small
document bundle. The project can compare suitable extraction approaches and evaluate
them against annotated examples.

The conveyancing use case also allows limited title reasoning. Once deed and property
facts are extracted, the system can help arrange the ownership history, identify missing
documents and flag possible inconsistencies such as name variations, plan-number
mismatches or extent differences. These checks are presented as support for lawyer
review, not as final legal conclusions.

\subsection{Software Engineering Component}
The software engineering component is the lawyer-facing workflow around the models:
matter creation, document upload, fact review, draft generation, editing and export.
The exact implementation choices can be decided later, but the project idea requires a
usable prototype that demonstrates the end-to-end flow from uploaded documents to
lawyer-reviewable drafts. Because legal documents are sensitive, privacy, access control
and auditability will be considered from the design stage.

\subsection{Evaluation and Objectives}
The system will be evaluated using measurable criteria:
\begin{itemize}[itemsep=0pt, topsep=0pt]
  \item Accuracy of document classification and key-field extraction
  \item Correctness of linked facts and simple title-chain support in the first use case
  \item Quality of missing-field and inconsistency warnings
  \item Lawyer assessment of generated drafts and source traceability
  \item Time saved compared with fully manual preparation
\end{itemize}

% =====================================================
\section{Datasets}

There is no ready-made public dataset for Sri Lankan legal document drafting workflows,
so building a small project dataset is part of the work. For the first use case, data
sources include anonymised historical deeds, title reports, survey plans, assessor
letters, ownership certificates and lawyer-approved drafting templates, collected with
permission from supporting practitioners. For general legal reference material, we can
refer to publicly available or subscription-based sources such as
LawLanka~\cite{lawlanka}, where appropriate.

On top of the raw documents, we will build a small annotated dataset that labels
entities such as parties, dates, deed numbers, boundaries, extents and plan references,
together with selected links between related documents. All documents will be anonymised
before use and handled under clear consent from the professionals who share them.

% =====================================================
\section{Similar Projects}

\noindent\textbf{Local.}
LawLanka, AI Pazz, Paralegal.lk and
LawCeylon~\cite{lawlanka,aipazz,paralegal,lawceylon}
are among the closest Sri Lankan systems.
Based on our preliminary review, these platforms mainly focus on legal search, case
law, legislation access and related research support.

\noindent\textbf{International.} Tools such as
Harvey AI, Spellbook, Lexis+ AI and
CoCounsel~\cite{harvey,spellbook,lexisai,cocounsel}
show that AI-assisted legal drafting is both viable and commercially valuable.
However, these are mostly built around other jurisdictions and legal workflows, so they
do not directly address the needs of Sri Lankan document drafting and review.

\noindent\textbf{Related research} Since Draftly is in the Sri Lankan
legal NLP domain, we will also review the previous work carried out under the SigmaLaw
project by Dr. Nisansa de Silva's research group~\cite{sigmalaw}. SigmaLaw has explored
legal information extraction, legal party identification, named entity recognition,
coreference resolution, semantic similarity, sentence embeddings, and legal information
retrieval, which are directly relevant to Draftly's broader goal of understanding legal
documents before drafting or review. The most relevant papers for our project include
``Party Identification of Legal Documents using Co-reference Resolution and Named Entity
Recognition''~\cite{sigmaparty}, ``Legal Party Extraction from Legal Opinion Texts
Using Recurrent Deep Neural Networks''~\cite{sigmarnn}, ``Legal Party Extraction from
Legal Opinion Text with Sequence to Sequence Learning''~\cite{sigmaseq2seq}, ``Context
Sensitive Verb Similarity Dataset for Legal Information Extraction''~\cite{sigmaverbsim},
``Learning Sentence Embeddings in the Legal Domain with Low Resource
Settings''~\cite{sigmasentemb}, and ``Legal Document Retrieval using Document Vector
Embeddings and Deep Learning''~\cite{sigmaretrieval}. We will use these works to
understand existing approaches for legal entity extraction, party linking, legal
semantic representation and information retrieval. Draftly differs by applying such
legal NLP ideas inside a lawyer-in-the-loop drafting and review product: the system
organises a client's uploaded legal documents, extracts and links facts with source
evidence, highlights missing or inconsistent information, and generates editable drafts
for professional review. Conveyancing and title documentation will be used as the first
practical workflow to validate this broader idea.

\noindent\textbf{How Draftly differs.} Draftly is positioned as a complementary system
for client-specific legal document drafting and review. Instead of only retrieving legal
sources, it helps organise uploaded matter documents, extract relevant facts, present
them for lawyer verification and generate editable drafts. Conveyancing and title
documentation are used as the first concrete workflow to prove this idea.

% =====================================================
\section{Scope and Team Responsibilities}

The project is scoped so that both a data-science and a software-engineering
contribution are clearly present and so that each team member owns a distinct part of
the system:
\begin{itemize}[itemsep=0pt, topsep=0pt]
  \item \textbf{Document understanding (Data Science)}: OCR support, entity extraction
        and comparison of extraction approaches.
  \item \textbf{Fact linking and checking (Data Science)}: linking related facts across
        documents and identifying missing or inconsistent information.
  \item \textbf{Draft generation (Data Science / SE)}: producing reviewable drafts from
        verified facts and approved templates.
  \item \textbf{Review workflow (Software Engineering)}: matter management, document
        upload, fact review, draft editing and export.
\end{itemize}

For the project idea, Draftly is broader than conveyancing, but the MVP is deliberately
limited to conveyancing and title-documentation workflows such as title reports,
pedigree generation, abstracts of title, ownership certificates and assessor letters.
After this first workflow is proven, the same product direction can be extended to
complaints, affidavits, petitions, legal notices and other Sri Lankan legal documents.

\begin{thebibliography}{00}

\bibitem{lawlanka}
LawLanka, ``Sri Lanka Acts \& Case Laws.'' [Online]. Available:
\url{https://www.lawlanka.com/}.

\bibitem{aipazz}
AI Pazz, ``AI-Powered Sri Lanka Legal Research Platform.'' [Online]. Available:
\url{https://www.aipazz.com/}.

\bibitem{paralegal}
Paralegal.lk, ``A search engine for Sri Lankan law.'' [Online]. Available:
\url{https://www.paralegal.lk/}.

\bibitem{lawceylon}
LawCeylon, ``Sri Lanka Acts, Laws \& Case Law Reports System.'' [Online].
Available: \url{https://lawceylon.lk/}.

\bibitem{harvey}
Harvey, ``AI software for legal and professional services.'' [Online]. Available:
\url{https://www.harvey.ai/}.

\bibitem{spellbook}
Spellbook, ``AI Contract Review \& Drafting.'' [Online]. Available:
\url{https://spellbook.com/}.

\bibitem{lexisai}
LexisNexis, ``Lexis+ with Protege: Legal AI Solution for Drafting \& Research.''
[Online]. Available:
\url{https://www.lexisnexis.com/en-us/products/lexis-plus-protege.page}.

\bibitem{cocounsel}
Thomson Reuters, ``CoCounsel Legal: AI Legal Assistant.'' [Online]. Available:
\url{https://legal.thomsonreuters.com/en/products/cocounsel-legal}.

\bibitem{sigmalaw}
Natural Language Processing Group, University of Moratuwa, ``SigmaLaw: Legal NLP
Research Project.'' [Online]. Available:
\url{https://nisansads.staff.uom.lk/Group/project-details.php?name=SigmaLaw}.

\bibitem{sigmaparty}
C. Samarawickrama, M. de Almeida, N. de Silva, G. Ratnayaka and A. S. Perera,
``Party Identification of Legal Documents using Co-reference Resolution and Named Entity
Recognition,'' in \emph{2020 IEEE 15th International Conference on Industrial and
Information Systems (ICIIS)}, 2020, pp. 494--499.
\url{https://doi.org/10.1109/ICIIS51140.2020.9342720}.

\bibitem{sigmarnn}
C. Samarawickrama, M. de Almeida, N. de Silva, G. Ratnayaka and A. S. Perera,
``Legal Party Extraction from Legal Opinion Texts Using Recurrent Deep Neural Networks,''
\emph{Journal of Data Intelligence}, vol. 3, no. 3, pp. 350--365, 2022.
\url{https://doi.org/10.26421/JDI3.3-4}.

\bibitem{sigmaseq2seq}
M. de Almeida, C. Samarawickrama, N. de Silva, G. Ratnayaka and A. S. Perera,
``Legal Party Extraction from Legal Opinion Text with Sequence to Sequence Learning,''
in \emph{2020 20th International Conference on Advances in ICT for Emerging Regions
(ICTer)}, 2020, pp. 143--148. \url{https://doi.org/10.1109/ICTer51097.2020.9325488}.

\bibitem{sigmaverbsim}
G. Ratnayaka, N. de Silva, A. S. Perera, G. Kavirathne, T. Ariyarathna and A. Wijesinghe,
``Context Sensitive Verb Similarity Dataset for Legal Information Extraction,''
\emph{Data}, vol. 7, no. 7, 2022. \url{https://doi.org/10.3390/data7070087}.

\bibitem{sigmasentemb}
S. Jayasinghe, L. Rambukkanage, A. Silva, N. de Silva, S. Perera and M. Perera,
``Learning Sentence Embeddings in the Legal Domain with Low Resource Settings,'' in
\emph{Proceedings of the 36th Pacific Asia Conference on Language, Information and
Computation (PACLIC)}, 2022, pp. 494--502. [Online]. Available:
\url{https://aclanthology.org/2022.paclic-1.55/}.

\bibitem{sigmaretrieval}
K. Sugathadasa, B. Ayesha, N. de Silva, A. S. Perera, V. Jayawardana, D. Lakmal and
M. Perera, ``Legal Document Retrieval using Document Vector Embeddings and Deep
Learning,'' in \emph{Science and Information Conference (SAI)}, 2018, pp. 160--175.
\url{https://doi.org/10.1007/978-3-030-01177-2_12}.

\end{thebibliography}

\end{document}
